% ============================================================================
% edicao.tex
% ----------------------------------------------------------------------------
% Este é o arquivo que você deve editar A CADA NOVO DOCUMENTO/RELATÓRIO.
% Aqui ficam as informações que mudam de um relatório para outro:
% título, subtítulo, nome do cliente, nome do projeto, equipe, e as opções
% de exibição do rodapé.
%
% Diferente de "config.tex" (que guarda dados fixos da empresa), este
% arquivo guarda dados VARIÁVEIS de cada edição/versão do documento.
% ============================================================================

% --------------------------------------------------------------
% ONDE FICA A NUMERAÇÃO DA PÁGINA
% -------------------------------------------------------------- 
% Comente a linha que não for usar e descomente a que for utilizar.
% \NumeracaoRodapetrue  -> mostra o número da página no RODAPÉ
% \NumeracaoRodapefalse -> mostra o número da página no CABEÇALHO

% \NumeracaoRodapetrue
\NumeracaoRodapefalse


% --------------------------------------------------------------
% TEXTO NO RODAPÉ
% -------------------------------------------------------------- 
% Comente a linha que não for usar e descomente a que for utilizar.
% \CNPJnoRodapetrue  -> mostra o CNPJ da empresa no rodapé
% \CNPJnoRodapefalse -> NÃO mostra o CNPJ no rodapé

% \CNPJnoRodapetrue
\CNPJnoRodapefalse


% --------------------------------------------------------------
% EMAIL EXIBIDO NO RODAPÉ
% -------------------------------------------------------------- 
% Descomente apenas a linha do setor responsável pelo documento.

\newcommand{\emailEmpresa}{contato@eletronjun.com.br}
% \newcommand{\emailEmpresa}{projetos@eletronjun.com.br}
% \newcommand{\emailEmpresa}{marketing@eletronjun.com.br}
% \newcommand{\emailEmpresa}{gdp@eletronjun.com.br}
% \newcommand{\emailEmpresa}{operacoes@eletronjun.com.br}
% \newcommand{\emailEmpresa}{presidencia@eletronjun.com.br}
% \newcommand{\emailEmpresa}{relacionamentos@eletronjun.com.br}
% \newcommand{\emailEmpresa}{comercial@eletronjun.com.br} 
% \newcommand{\emailEmpresa}{comissaoeleitoral@eletronjun.com.br}


% --------------------------------------------------------------
% NOME DO PROJETO, CLIENTE, GERENTE E EQUIPE TÉCNICA
% -------------------------------------------------------------- 
\newcommand{\nomeProjeto}{Sistema de Monitoramento Ambiental}   % nome do projeto entregue
\newcommand{\nomeCliente}{Eletronjun}   % nome da empresa/cliente atendido

\newcommand{\gerenteProjeto}{João Gabryel da Silva Costa}    % responsável pelo projeto
\newcommand{\equipeTecnica}{Laís Helena, Ana Clara Patrício, Kelyton de Lucas, Pedro Henrick} % nomes da equipe técnica (separados por vírgula)

% --------------------------------------------------------------
% TÍTULO E SUBTÍTULO DO DOCUMENTO
% -------------------------------------------------------------- 
\newcommand{\tituloDocumento}{Documentação Técnica — Projeto Leitor de Temperatura}
\newcommand{\subtituloDocumento}{Processo Seletivo Eletronjun 2026.2}
